\section{Orientación trítica, prolongación compatible y doble proyección}
\label{sec:trit-projection}

El carácter generado de las coordenadas reales se expresa en la relación
entre profundidad aritmética y evaluación. El dato primitivo es una historia
admisible de operaciones con sus residuos, cocientes y orientaciones. Una
lectura arquimediana asocia a sus prefijos intervalos cada vez más estrechos;
una lectura incidencial conserva relaciones de frontera del mismo estado.
La doble proyección del corpus \cite{hmtI,HMT-IX} reúne esas dos operaciones
sin convertir una evaluación escalar en la totalidad de la historia. El
artículo utiliza ese antecedente para estudiar las realizaciones posteriores
del registro dodecafásico.

\subsection{Cilindros posicionales y determinación de la coordenada real}

Para un bloque ternario \(b=(b_1,\ldots,b_6)\), con representantes
\(b_j\in\{0,1,2\}\) de \(\mathbb F_3\), la evaluación entera
\(\nu(b)=\sum_{j=1}^6b_j3^{6-j}\) pertenece a
\(\{0,\ldots,728\}\). Si \(b_{n+1}\) es el bloque admisible siguiente,
la prolongación posicional satisface
\[
 N_{n+1}=729N_n+\nu(b_{n+1}),\qquad
 a_n=\frac{N_n}{729^n},\qquad b_n^+=\frac{N_n+1}{729^n}.
\]
Aquí \(N_0\) es la parte entera que determina el lector regional. La letra
\(b_n^+\) designa un extremo real y se distingue del bloque ternario \(b\).
El cilindro de representación es \([a_n,b_n^+)\); su clausura se utiliza
en el argumento de existencia del límite. El acarreo conserva el paso de
un prefijo al anterior mediante división euclídea.

\begin{proposition}[Límite de los centros cilíndricos]
Para una prolongación compatible, las clausuras de los cilindros están
encajadas y determinan una única coordenada \(x\). Si
\(c_n=(a_n+b_n^+)/2\), entonces
\[
 x=\lim_{n\to\infty}a_n=\lim_{n\to\infty}b_n^+
 =\lim_{n\to\infty}c_n,\qquad
 |x-c_n|\leq\frac1{2\,729^n}.
\]
\end{proposition}
\begin{proof}
La desigualdad \(0\leq\nu(b_{n+1})<729\) sitúa el cilindro sucesor en
la clausura del predecesor. Sus extremos son racionales, su anchura es
\(729^{-n}\) y permanecen dentro del primer intervalo cerrado acotado.
El encajamiento y la anchura decreciente producen una única clase de
Cauchy de extremos. En la realización arquimediana, esa clase tiene el
representante \(x\); la distancia al punto medio es a lo sumo la mitad
de la anchura. La convención de frontera asigna una escritura cuando dos
expansiones posicionales representan el mismo límite.
\end{proof}

La media aritmética interviene, por tanto, en una evaluación precisa: el
centro de cada intervalo que acota la coordenada. El límite conserva la
exactitud del valor aunque la anchura desaparezca. La representación
enriquecida conserva adicionalmente la sucesión de prefijos y su incidencia.
La proyección numérica y la proyección excepcional son dos aplicaciones
con codominios diferentes; una conserva el valor límite y la otra transporta
las relaciones que habilitan la construcción reticular. El paso a
\(\RR\) se sitúa así en la realización del lector, después de la
generación aritmética de los datos que evalúa.

Para hacer explícita la compatibilidad, sean \(p_{n+1,n}\) las restricciones
del estado enriquecido y \(x_n=p_{n+1,n}(x_{n+1})\). Si \(I_n(x_n)\)
es el cilindro numérico y \(\mathcal B_n(x_n)\) la incidencia de frontera,
los dos transportes conservan simultáneamente
\[
 \overline{I_{n+1}(x_{n+1})}\subseteq\overline{I_n(x_n)},\qquad
 r_{n+1,n}\mathcal B_{n+1}(x_{n+1})=\mathcal B_n(x_n).
\]
La segunda igualdad utiliza los mapas incidenciales del corpus; la
proposición~\ref{prop:lf-incidence-naturality} explicita su realización
mediante marcos y códigos transportados. Sus cinco
construcciones consustanciales ---espectral, solenoidal, gauge, cohomológica
y determinantal--- pertenecen a ese mismo sistema compatible. El estudio
geométrico posterior conserva esta unidad de origen al especificar cuál
de sus observaciones utiliza en cada teorema.

\subsection{Trayectorias orientadas y temporalidad de la actualización}

Sea \(x\xrightarrow{\rm adm}y\) la relación de admisibilidad del TPK en
el espacio enriquecido \(\widehat X\). Una trayectoria parametrizada por
un conjunto discretamente ordenado \(\mathcal T\) es
\[
 \gamma:\mathcal T\longrightarrow\widehat X,
 \qquad\gamma(t)\xrightarrow{\rm adm}\gamma(t^+).
\]
La relación de compatibilidad determina las transiciones; la parametrización
permite ordenarlas. En el espacio \(E=\mathcal T\times\widehat X\),
con \(p(t,x)=t\), la trayectoria determina la sección
\(\widetilde\gamma(t)=(t,\gamma(t))\), que satisface
\(p\circ\widetilde\gamma=\id_{\mathcal T}\).
Esta precisión conserva la diferencia entre un estado, su trayectoria y
la coordenada utilizada para describir la evolución.

La firma trítica \(\tau\in\{+1,0,-1\}\) tipa los generadores mediante
\[
 J_\tau^2=-\tau I.
\]
En el lector temporal del corpus, \(+1\) representa retorno y memoria;
\(0\), actualización de umbral; y \(-1\), apertura de prolongaciones
conjugadas. Su interpretación como pasado, presente y futuro se refiere
a las funciones del antecedente conservado, la transición efectiva y la
continuación admisible. Los tres tiempos verbales adquieren así un
contenido operatorio sobre historias.

Si \(\mathcal P_n(x)\) es el conjunto de continuaciones admisibles de
longitud \(n\), la supresión del último paso define
\(r_n:\mathcal P_{n+1}(x)\to\mathcal P_n(x)\). Una continuación
ilimitada es una colección \((\gamma_n)\) con
\(r_n(\gamma_{n+1})=\gamma_n\). La restricción reconstruye antecedentes;
la actualización \(U_t\) actúa en la transición; la prolongación mantiene
la frontera del prefijo. La existencia de estas colecciones se utiliza en
los dominios supervivientes fijados por el teorema de prolongación del
corpus.

Sobre una firma finita \(\sigma=(\tau_1,\ldots,\tau_n)\), la inversión
orientada es
\[
 C_{\rm rev}\sigma=(-\tau_n,\ldots,-\tau_1),
 \qquad C_{\rm rev}^2\sigma=\sigma.
\]
La trayectoria conjugada conserva también las transformaciones de hoja,
posición y frontera de su transporte. Por ello la bidireccionalidad se
expresa en la historia completa, mientras esta fórmula expresa su firma.
El retorno de fase de \(\Gamma_9\) admite un incremento de memoria:
estados con fase idéntica pueden ocupar profundidades diferentes.

\subsection{Álgebras cuadráticas y actividad del umbral parabólico}

En un régimen fijo, la conjugación de \(aI+bJ_\tau\) cambia el signo
de \(b\), y su norma algebraica es
\[
 (aI+bJ_\tau)(aI-bJ_\tau)=(a^2+\tau b^2)I.
\]
Los tipos elíptico, parabólico e hiperbólico corresponden a las tres
relaciones cuadráticas. Sus grupos locales se representan por
\[
 e^{tJ_{+1}}=\cos t\,I+\sin t\,J_{+1},\quad
 e^{tJ_0}=I+tJ_0,\quad
 e^{tJ_{-1}}=\cosh t\,I+\sinh t\,J_{-1}.
\]
La prueba separa términos pares e impares en la serie exponencial; en el
tipo nilpotente sólo permanecen los grados cero y uno. Cada colección
satisface la ley de composición en su parámetro y la inversión
\(e^{-tJ_\tau}=(e^{tJ_\tau})^{-1}\).

El valor trítico cero permite, en particular, un generador distinto de
cero con cuadrado nulo. Para
\(N=\left(\begin{smallmatrix}0&1\\0&0\end{smallmatrix}\right)\),
\[
 e^{tN}(q,p)=(q+tp,p).
\]
El umbral tiene una acción lineal efectiva sobre su fibra. Su parámetro,
su desplazamiento y la memoria de la transición permanecen determinados
aunque la firma de régimen sea cero. Éste es el contenido preciso que
permite desarrollar narrativamente el presente como actualización activa.

La forma elíptica positiva tiene niveles circulares en coordenadas
ortonormales y elípticos tras una transformación lineal; la forma hiperbólica
posee dos direcciones isotrópicas y niveles hiperbólicos; el tipo parabólico
posee una forma degenerada y transporte unipotente. Una geometría esférica
requiere además el espacio homogéneo o la métrica que la realiza. La
clasificación anterior fija los generadores que pueden intervenir en esa
realización y mantiene separados el tipo algebraico y la curvatura de un
espacio posterior.

El calendario de modos concreta la conexión con la autoescala. La regla
\(+1:m\mapsto\Sigma\), \(-1:m\mapsto\Pi\), \(0:m\mapsto m\)
aplicada al bloque cíclico \((+1,-1,0)\) produce
\((\Sigma,\Pi,\Pi)\). Contar sus tres transiciones, en el orden
\((\Sigma,\Pi)\), determina
\begin{equation}
 F_{\rm av}=\begin{pmatrix}0&1\\1&1\end{pmatrix}.
 \label{eq:geo-fav}
\end{equation}
Su rayo propio positivo, normalizado como \((1,x)^\mathsf T\), satisface
\(x=\lambda\), \(1+x=\lambda x\), y por ello
\(x^2-x-1=0\). La raíz positiva es la coordenada de autoescala expresada
como \(\varphi\). La matriz nace del calendario de operaciones; la
expresión radical y las identidades exponenciales siguientes son sus
realizaciones posteriores.

% Articulación posterior a la generación de pi, phi y e.
% Contenido acumulativo; no reemplaza el desarrollo previo del TRIT.
\subsection{Realizaciones exponenciales de las coordenadas generadas}
\label{geo-euler-trit-articulacion}

Las coordenadas regionales ya determinadas permiten reconocer las
realizaciones concretas de la colección cuadrática del TRIT. El orden es
sustantivo: la construcción de las regiones precede a la evaluación de una
exponencial en los parámetros \(\pi\) o \(2\log\varphi\). Estas expresiones
identifican la función de las coordenadas en operadores construidos desde
APP y TPK; no intervienen en la selección de los prefijos que las producen.

\subsubsection{Ciclo ternario, transporte nilpotente e incidencia de hojas}

La multiplicación \(a(r)=4r\) en \(\mathbb Z/9\mathbb Z\) tiene orden
\(3\). Sobre las unidades determina las órbitas \(\{1,4,7\}\) y
\(\{2,8,5\}\). En el módulo de aumento de cualquiera de ellas, formado por
las combinaciones enteras cuyos coeficientes suman cero, la base
\((e_r-e_{4r},e_{4r}-e_{7r})\) da
\[
 C=\begin{pmatrix}0&-1\\1&-1\end{pmatrix},
 \qquad C^2+C+I=0.
\]
La forma restringida tiene matriz de Gram
\(\bigl(\begin{smallmatrix}2&-1\\-1&2\end{smallmatrix}\bigr)\).
En las evaluaciones residuales módulo \(9\), la acción distingue las hojas:
\(S(ax,ay)=aS(x,y)\), mientras \(P(ax,ay)=a^{-1}P(x,y)\).
La realización cíclica conserva, por tanto, una orientación contragrediente
entre suma y producto. Los cocientes del levantamiento entero se conservan
en el registro aritmético correspondiente.

El transporte parabólico elemental se representa mediante
\[
 N=\begin{pmatrix}0&1\\0&0\end{pmatrix}.
\]
La incidencia areal--volumétrica ya construida en \eqref{eq:geo-fav} proporciona
\[
 F_{\rm av}=\begin{pmatrix}0&1\\1&1\end{pmatrix},
 \qquad A=F_{\rm av}^2=\begin{pmatrix}1&1\\1&2\end{pmatrix}.
\]
Los tres operadores actúan en espacios reales declarados: el módulo de
aumento tensorizado con \(\mathbb R\), el plano del transporte nilpotente y
el plano de incidencia modal. Compartir dimensión matricial no identifica
estos espacios ni sus coordenadas.

\begin{proposition}[Generadores concretos de los tres tipos]
\label{geo-euler-trit-generadores}
Los operadores
\[
 J_{\mathrm e}=\frac{2C+I}{\sqrt3},\qquad
 J_{\mathrm p}=N,\qquad
 J_{\mathrm h}=\frac{2A-3I}{\sqrt5}
\]
satisfacen, respectivamente,
\[
 J_{\mathrm e}^{\,2}=-I,\qquad
 J_{\mathrm p}^{\,2}=0,\qquad
 J_{\mathrm h}^{\,2}=I.
\]
\end{proposition}
\begin{proof}
La relación de \(C\) implica \((2C+I)^2=-3I\).
La multiplicación directa da \(N^2=0\).
Finalmente, \(A\) tiene traza \(3\) y determinante \(1\), por lo que
\(A^2-3A+I=0\) y \((2A-3I)^2=5I\).
\end{proof}

\begin{proposition}[Tres evaluaciones de la identidad exponencial]
\label{geo-euler-trit-evaluaciones}
En las realizaciones anteriores,
\begin{equation}
 \boxed{
 \exp(\pi J_{\mathrm e})+I=0,\qquad
 \exp(J_{\mathrm p})=I+J_{\mathrm p},\qquad
 \exp(2\log\varphi\,J_{\mathrm h})=F_{\rm av}^{\,2}.}
 \label{geo-euler-trit-tres-evaluaciones}
\end{equation}
\end{proposition}
\begin{proof}
La especialización elíptica de la exponencial trítica da
\(\cos\pi\,I+\sin\pi\,J_{\mathrm e}=-I\).
La nilpotencia elimina exactamente los términos de grado al menos \(2\)
en la segunda expresión. Para la tercera, \(\varphi^2=\varphi+1\) implica
\[
 \cosh(2\log\varphi)=\frac32,\qquad
 \sinh(2\log\varphi)=\frac{\sqrt5}{2}.
\]
Por tanto, la exponencial vale
\(\frac32I+\frac{\sqrt5}{2}J_{\mathrm h}=A\).
\end{proof}

La primera igualdad representa la media vuelta elíptica; la vuelta completa
satisface \(\exp(2\pi J_{\mathrm e})=I\). La segunda mantiene un término
lineal no nulo: el régimen parabólico expresa transporte, no ausencia de
evolución. La tercera realiza un avance hiperbólico unimodular. Sus
autovalores son \(\varphi^2\) y \(\varphi^{-2}\); una dirección se expande y
la otra se contrae conservando el determinante. Las tres evaluaciones
emplean parámetros distintos de una colección exponencial común.

En esa colección, \(e\) interviene mediante la evaluación escalar
\(\exp(I)=eI\) y la ley \(\exp(sI)\exp(tI)=\exp((s+t)I)\).
Su papel no se restringe al tipo parabólico. La coordenada \(e\), obtenida
previamente de su región, es reconocida aquí como la evaluación unitaria
de la exponencial; esta identidad no sustituye su genealogía coinductiva.

\subsubsection{Resolución polinómica de los regímenes}

Las tres realizaciones admiten una presentación compacta sobre la suma
directa de sus espacios. Defínanse
\[
 \mathbb J=J_{\mathrm e}\oplus J_{\mathrm p}\oplus J_{\mathrm h},
 \qquad Q=\mathbb J^2.
\]
Las relaciones anteriores implican
\[
 \mathbb J^6=\mathbb J^2,\qquad Q^3=Q.
\]
El uso de \(Q\) mantiene separado este cuadrado operatorial del registro
dodecafásico \(K\).

\begin{proposition}[Proyectores polinómicos de régimen]
\label{geo-euler-trit-proyectores}
Las aplicaciones
\[
 p_{\mathrm e}=\frac{Q^2-Q}{2},\qquad
 p_{\mathrm p}=I-Q^2,\qquad
 p_{\mathrm h}=\frac{Q^2+Q}{2}
\]
son idempotentes mutuamente ortogonales, suman \(I\) y recuperan los tres
sumandos de la representación.
\end{proposition}
\begin{proof}
Los tres polinomios toman los valores \((1,0,0)\), \((0,1,0)\) y
\((0,0,1)\) en los puntos \(-1,0,+1\), respectivamente. Las diferencias
entre sus cuadrados y ellos mismos, así como sus productos cruzados, son
múltiplos de \(X^3-X\). Sustituir \(Q\), que anula ese polinomio, demuestra
las identidades. En los tres bloques, \(Q\) actúa como \(-I,0,I\).
\end{proof}

Se obtiene así
\begin{align}
 \exp(t\mathbb J)
 ={}&p_{\mathrm e}(\cos t\,I+\sin t\,\mathbb J)
   +p_{\mathrm p}(I+t\mathbb J)\notag\\
   &+p_{\mathrm h}(\cosh t\,I+\sinh t\,\mathbb J).
 \label{geo-euler-trit-exponencial-total}
\end{align}
Cada sumando conserva su relación cuadrática y la conjugación invierte
su evolución. La representación total reúne los tres regímenes, pero no
define por sí sola un operador de transición entre ellos. Tal transición
requiere los mapas de transporte del TPK con sus dominios, orientación y
memoria. La articulación exponencial permite reconocer sus tipos sin
reemplazar esa dinámica por una identificación de las fibras.

% PROCEDENCIA FOCAL (no impresa)
% Resultado recuperado: tres generadores y tres evaluaciones exponenciales.
% Propietario completo:
% BASE/manuscrito/sections/hmt/09b_algebras_cuadraticas.tex:289-516.
% Antecedente causal de C y acción contragrediente:
% BASE/manuscrito/sections/hmt/03d_esqueleto_ciclotomico_app.tex:13-47,145-246.
% Los tres tipos y la orientación están en:
% BASE/manuscrito/sections/hmt/02_trit_geometrias.tex:83-196.
% Resolución polinómica recuperada:
% output/INVESTIGACION_HOLONOMIA_NONADICA_CRISTAL_APERIODICO_20260903/
% ALGEBRA_HOLONOMICA_UNIVERSAL_Y_EULER_TRITICO.md, secciones 8-11 y 22.
% Los cuatro propietarios se leyeron completos. BASE designa:
% /Users/ruben/Documents/HMT2/
% EL_CIERRE_HOLOGRAFICO_DEL_INFINITO_SUCESOR_102_CAPITULOS_2026-08-28_EN_TRABAJO/
% No se incorpora la identidad P9(alpha)=delta del documento especializado:
% esa identificación no pertenece a este bloque.
% Tampoco se identifica una dirección nilpotente con la vacancia electrónica,
% ni se presume una representación global del transporte por la suma directa.
% COMPROBACIONES: Python estándar, enteros y Fraction.
% C^2+C+I=0; (2C+I)^2=-3I; N^2=0; (2Fav^2-3I)^2=5I.
% 81 pares de residuos verifican las dos acciones contragredientes.
% Los 3 proyectores se verificaron en Q=-1,0,+1; prueba general en texto.
% 6 labels únicos y entornos equilibrados. Sin compilación.

